% !TeX program = lualatex
\makeatletter
\def\input@path{{00/plantilla/fisica/}}
\makeatother
\documentclass[carta,firma]{lafrox}
\usepackage{etoolbox}
\captionsetup{hypcap=false}
\setcounter{tocdepth}{2}
\setcounter{secnumdepth}{2}
% Índice detallado compacto: conserva tamaño y títulos de la plantilla.
\titlecontents{chapter}[0pt]
  {\addvspace{7pt}\intersemibold\color{lafroxTinta}}
  {\textcolor{lafroxFuerte}{\thecontentslabel}\enspace}{}
  {\normalfont\color{lafroxGris}\hfill\contentspage}
\titlecontents{section}[1.5em]
  {\addvspace{1pt}\fontsize{10pt}{12pt}\selectfont\color{lafroxTinta}}
  {\textcolor{lafroxGris}{\thecontentslabel}\enspace}{}
  {\normalfont\color{lafroxGris}\hfill\contentspage}
\titlecontents{subsection}[3.4em]
  {\fontsize{9.5pt}{11pt}\selectfont\color{lafroxGris}}
  {\thecontentslabel\enspace}{}
  {\hfill\contentspage}
\newcommand{\datosdocumentocinco}[2]{
\datoslafrox{documento={Propuesta Técnica},subtitulo={#1},alcance={Modelo y gestión de datos},formulario={Formulario T-7},fecha={02 de octubre de 2026},version={1.0},nomenclatura={#2}}
\hypersetup{pdftitle={#1}}
\renewcommand{\LFXestadodoc}{Propuesta de diseño}
\patchcmd{\portadalafrox}{\LFXcampoportada{Versión}{\LFXversion}}{\LFXcampoportada{Documento}{5}}{}{}
}
